%% ma: L11-B24-0003
%% lop: 11
%% chuong: 7
%% bai: 24
%% dang: 11.24.3
%% loai: TN4
%% muc: TH
%% dap_an: "B"
%% nguon: "(NBV)-ĐỀ SỐ 11-MỖI NGÀY 1 ĐỀ THI 2025.pdf (D11, MỖI NGÀY 1 ĐỀ - Đề số 11), Phần I Câu 4"
%% kien_thuc: "Bài 24 (góc giữa đường thẳng và mặt phẳng, phép chiếu vuông góc), Bài 23 (đường thẳng vuông góc với mặt phẳng), Bài 11 (lăng trụ đứng, hình vuông)"
%% trang_thai: cho-duyet
%% ly_do: "Dạng toán mới đề xuất, chờ giáo viên duyệt."
%% ghi_chu: "Đáp án gốc B đúng (sympy: tan = căn 3). Đề gốc không có hình."
\begin{ex}
Cho hình lăng trụ tứ giác đều $ABCD.A'B'C'D'$ có $AB=a$, $AA'=a\sqrt6$. Góc giữa $AC'$ và mặt phẳng $(A'B'C'D')$ bằng
\choice{$30^\circ$}{\True $60^\circ$}{$45^\circ$}{$90^\circ$}
\loigiai{Vì $AA'\perp(A'B'C'D')$ nên $A'C'$ là hình chiếu vuông góc của $AC'$ lên $(A'B'C'D')$. Do đó góc cần tìm là $\widehat{AC'A'}$.

$A'C'=a\sqrt2$ (đường chéo hình vuông cạnh $a$). Trong tam giác $AA'C'$ vuông tại $A'$: $\tan\widehat{AC'A'}=\dfrac{AA'}{A'C'}=\dfrac{a\sqrt6}{a\sqrt2}=\sqrt3$, suy ra $\widehat{AC'A'}=60^\circ$.}
\end{ex}
